\section{Introduction}

Long-horizon time-series forecasting papers usually treat the model architecture as the primary design object. Attention sparsity, temporal decomposition, patching, channel mixing, and embedding design receive direct experimental attention, while the feed-forward activation is often inherited from a reference implementation. This convention is convenient, but it hides a real modeling choice. The activation function determines whether intermediate representations are sign-preserving or rectified, bounded or unbounded, smooth or piecewise linear, and whether large excursions are amplified or saturated. These properties are especially relevant in forecasting, where the same architecture is evaluated across horizons, frequencies, and datasets with different volatility and seasonal structure.

This paper asks a deliberately narrow question: if the Informer architecture and training protocol are held fixed, how much does the activation configuration matter? We choose Informer \citep{zhou2021informer} as the controlled substrate because it is a canonical long-sequence forecasting transformer and its public scripts define horizon-specific settings for the ETT benchmarks. We do not propose a new attention mechanism, loss function, or backbone. Instead, we expose the activation family and its placement as explicit experimental variables while keeping attention, embedding, depth, width, optimizer, data split, batch size, and epoch budget fixed inside each dataset-horizon cell.

The resulting answer is positive, but not in the form of a universal activation rule. In the completed Informer grid, non-GELU activations are best in 34 of 36 dataset-horizon cells. Bounded and tanh-family activations are the strongest static choices: \texttt{tanh\_sin001\_all}, \texttt{tanh\_all}, and \texttt{softsign\_all} have the best mean ranks and together win 23 cells. At the same time, the best configuration changes across datasets and horizons. Solar3 remains GELU-best at both tested horizons, output \texttt{tanh} wins several ETT/ETTm cells but is poor on average, and ReLU wins only one cell. Thus the main result is not that one activation replaces GELU everywhere. It is that activation choice is a first-order design variable under a fixed forecasting protocol.

\Cref{fig:phase4_best_gain_heatmap} summarizes this pattern. The large gains concentrate in ETTh2 and ETTm2, while Solar gains are smaller and site-specific. This heterogeneity motivates a second question: can input-window factors such as volatility, jump intensity, turbulence, or spectral concentration select the best activation without retraining? Our current evidence says no. Factor bins reveal regime dependence, but leave-one-seed-out oracle evidence is weak and localized. We therefore use factors as diagnostics for targeted follow-up studies, not as a deployed zero-shot router.

We also test whether the Informer ranking transfers to other forecasting architectures. PatchTST \citep{nie2023patchtst} shows partial transfer: non-GELU activations win selected ETTh2 and Solar cells, but GELU remains the best static choice in most tested cells, and output \texttt{tanh} loses every paired comparison against the matching linear head. iTransformer \citep{liu2023itransformer} is an even stricter stress test: in a one-seed FFN-only pilot, GELU remains best in five of six cells, while a later small placement probe suggests conditional output-tanh gains on selected cells. These results support an architecture-dependent story and prevent overclaiming direct transfer.

The contributions are:
\begin{enumerate}
    \item We implement an activation-only Informer evaluation interface that changes hidden activation family and placement while holding the rest of the architecture and training protocol fixed.
    \item We report a complete three-seed grid over 36 dataset-horizon cells, nine activation configurations, and 972 matched seed-runs, showing that non-GELU activations are best in 34/36 cells.
    \item We identify bounded and tanh-family activations as strong static Informer choices while showing that no single activation or placement universally dominates.
    \item We separate diagnostic regime evidence from deployment claims: factor bins expose conditional behavior, but current leave-one-seed-out results do not support a zero-shot activation router.
    \item We include PatchTST and iTransformer stress tests to show that activation sensitivity is real but mediated by architecture and placement.
\end{enumerate}

% Generated by figures/gen_phase4_paper_assets.py
\begin{figure*}[t]
    \centering
    \includegraphics[width=\textwidth]{figures/fig_phase4_best_gain_heatmap.pdf}
    \caption{Best activation gain over the same-protocol GELU row in each completed dataset-horizon cell. The full grid covers 36 dataset-horizon cells, nine activation configurations, and three seeds. Most cells prefer a non-GELU activation, but Solar3 remains GELU-best at both horizons.}
    \label{fig:phase4_best_gain_heatmap}
\end{figure*}

\begin{figure}[t]
    \centering
    \includegraphics[width=\linewidth]{figures/fig_phase4_static_ranking.pdf}
    \caption{Static activation ranking across the full Phase-4 grid. Bounded and tanh-family activations rank best on average, but the best-cell counts are spread across several configurations, arguing against a universal replacement activation.}
    \label{fig:phase4_static_ranking}
\end{figure}

